%% schriftsatz.latex — Anpassbares Pandoc-Template für familienrechtliche Schriftsätze
%%
%% Pandoc-Variablen (via -V oder YAML-Frontmatter):
%%   mainfont      Schriftart             (Standard: Arial)
%%   fontsize      Schriftgröße           (Standard: 11pt)
%%   papersize     Papierformat           (Standard: a4)
%%   geometry      Seitenränder           (Standard: margin=2.5cm)
%%   linestretch   Zeilenabstand          (Standard: 1.2)
%%   lang          Sprache                (Standard: de-DE)
%%   briefkopf     Vorbereiteter LaTeX-Briefkopf (von generate-pdf.py erzeugt)
%%
%% ─────────────────────────────────────────────────────────────────────────────

\documentclass[
  11pt,
  a4paper,
]{article}

%% ── Sprache ──────────────────────────────────────────────────────────────────
\usepackage[ngerman]{babel}

%% ── Fonts (XeLaTeX) ──────────────────────────────────────────────────────────
\usepackage{fontspec}
\setmainfont{Arial}[
  BoldFont       = Arial Bold,
  ItalicFont     = Arial Italic,
  BoldItalicFont = Arial Bold Italic,
  Ligatures      = TeX,
]
\usepackage{microtype}

%% ── Seitenlayout ─────────────────────────────────────────────────────────────
\usepackage[left=3cm,right=2cm,top=2.5cm,bottom=2.5cm]{geometry}

%% ── Zeilenabstand & Absätze ──────────────────────────────────────────────────
\usepackage{setspace}
\setstretch{1.3}
\usepackage{parskip}
\setlength{\parskip}{6pt}
\setlength{\parindent}{0pt}

%% ── Überschriften ────────────────────────────────────────────────────────────
\usepackage{titlesec}
\titleformat{\section}{\normalfont\large\bfseries}{}{0em}{}
\titleformat{\subsection}{\normalfont\normalsize\bfseries}{}{0em}{}
\titleformat{\subsubsection}{\normalfont\normalsize\bfseries}{}{0em}{}
\titlespacing{\section}{0pt}{14pt}{6pt}
\titlespacing{\subsection}{0pt}{10pt}{4pt}
%% Keine automatische Abschnittsnummerierung
\setcounter{secnumdepth}{0}

%% ── Seitenzahlen ─────────────────────────────────────────────────────────────
\usepackage{fancyhdr}
\usepackage{lastpage}
\pagestyle{fancy}
\fancyhf{}
\fancyfoot[C]{\footnotesize Seite~\thepage~von~\pageref{LastPage}}
\renewcommand{\headrulewidth}{0pt}
\renewcommand{\footrulewidth}{0pt}

%% ── Tabellen ─────────────────────────────────────────────────────────────────
\usepackage{booktabs}
\usepackage{longtable}
\usepackage{array}
\usepackage{calc}
\usepackage{caption}
\captionsetup{font=small, labelfont=bf}
%% Pandoc 3.x setzt \LTcaptype auf "none" — LaTeX braucht diesen Counter:
\newcounter{none}

%% ── Listen ───────────────────────────────────────────────────────────────────
\usepackage{enumitem}
\setlist{noitemsep, topsep=4pt}

%% ── Sonstiges ────────────────────────────────────────────────────────────────
\usepackage{xcolor}
\usepackage[normalem]{ulem}
\usepackage[hidelinks]{hyperref}
\usepackage{needspace}

%% ── Pandoc-Kompatibilität ────────────────────────────────────────────────────
\providecommand{\tightlist}{%
  \setlength{\itemsep}{0pt}\setlength{\parskip}{0pt}}

%% ─────────────────────────────────────────────────────────────────────────────
\begin{document}

%% Briefkopf — von generate-pdf.py aus Markdown-Zeilen vor dem ersten --- erzeugt.
%% Zeilen mit "den \d" werden rechtsbündig gesetzt, Betreff fett+unterstrichen,
%% Aktenzeichen fett, alle anderen Zeilen linksbündig mit explizitem Zeilenumbruch.

\section{Nur mündlich --- Az. 7 F
88/25}\label{nur-muxfcndlich-az.-7-f-8825}

\subsection{Punkte für die mündliche
Verhandlung}\label{punkte-fuxfcr-die-muxfcndliche-verhandlung}

\subsubsection{Spontanes Lob der KiTa-Leiterin nach dem
Entwicklungsgespräch}\label{spontanes-lob-der-kita-leiterin-nach-dem-entwicklungsgespruxe4ch}

\textbf{Kontext:} Die KiTa-Leiterin hat AG nach dem formellen
Entwicklungsgespräch am 25.02.2026 spontan und ohne Anlass gelobt. Das
Zitat ist zu persönlich und klingt im Schriftsatz aufgesetzt; es würde
zudem der Gegenseite ermöglichen, die KiTa als befangen darzustellen.

\textbf{Kernaussage:} Selbst die unmittelbare, unvorbereitete Reaktion
der KiTa-Leiterin spricht für die positive Entwicklung Finns und die
Qualität der Betreuung durch den Antragsgegner.

\textbf{Redevorschlag:}

\begin{center}\rule{0.5\linewidth}{0.5pt}\end{center}

\subsubsection{Reaktion auf Mitbringen der Schwiegermutter bei
Übergabe}\label{reaktion-auf-mitbringen-der-schwiegermutter-bei-uxfcbergabe}

\textbf{Kontext:} AS hat am 28.02.2026 erstmals ihre Mutter zur Übergabe
mitgebracht. Das klingt im Schriftsatz kleinlich, könnte aber auf
eskalierende Dynamik hinweisen. Mündlich lässt es sich als ruhige
Beobachtung einbringen, falls die Richterin nach der
Kommunikationssituation fragt.

\textbf{Kernaussage:} Die Übergaben funktionierten bislang reibungslos;
AG möchte das nicht eskalieren, beobachtet aber eine Veränderung seit
Einreichung des Antrags.

\textbf{Redevorschlag:}

\begin{center}\rule{0.5\linewidth}{0.5pt}\end{center}

\subsection{Gründe für „nur
mündlich''}\label{gruxfcnde-fuxfcr-nur-muxfcndlich}

Punkte landen hier wenn: - Sie emotional sind und im Schriftsatz
übertrieben wirken würden - Sie der Gegenseite Angriffsfläche bieten,
aber mündlich kontextualisiert werden können - Sie Details enthalten,
die nur auf Nachfrage relevant sind - Sie den Schriftsatz aufblähen,
ohne Beweiswert hinzuzufügen - Sie nur als Reaktion auf erwartbare
Fragen/Argumente vorbereitet werden

\end{document}
